
\subsubsection{Experimental Design}

Experiments are designed to isolate the effect of attention structure on attribution method performance. The key challenge is controlling for confounding factors---model size, depth, training procedure---that could obscure architecture-specific effects.

\textbf{Architecture Selection.} BERT-base-uncased (110M parameters, 12 layers) is compared with GPT-2 (124M parameters, 12 layers). Both models have identical layer counts, similar parameter counts, use the same hidden dimension (768), and were pretrained on large text corpora. The primary architectural difference is attention structure: BERT uses bidirectional attention, while GPT-2 uses causal attention.

\textbf{Task Selection.} SST-2 binary sentiment classification from the GLUE benchmark is used. Both architectures can perform it natively (BERT: [CLS] classification, GPT-2: last-token classification). 5\% random label noise is injected to create ground-truth mislabeled examples (3,367 mislabeled out of 67,349 training examples).

\textbf{Training Protocol.} Both models are fine-tuned with identical hyperparameters: AdamW optimizer, learning rate 2e-5, 3 epochs, batch size 32.

\subsubsection{Attribution Methods}

Three representative methods are evaluated:

\begin{itemize}
\item \textbf{TRAK} \citep{park2023trak}: Projects gradients to random subspace, performs linear regression. Projection dimensions: {64, 256, 1024}.
\item \textbf{EK-FAC} \citep{grosse2023studying}: Kronecker-factored Hessian approximation. Projection dimensions: {64, 256, 1024}.
\item \textbf{TracIn} \citep{pruthi2020estimating}: Gradient dot-products at checkpoints. Checkpoints: {1, 2, 3}.
\end{itemize}

\subsubsection{Causal Mechanism Verification}

The hypothesized causal chain is verified through four sub-hypotheses:
\begin{itemize}
\item \textbf{H-M1:} Attention sparsity differs >90\% between architectures
\item \textbf{H-M2:} Hessian eigenvalue differs >10\% between architectures
\item \textbf{H-M3:} All methods show >10\% architecture effect in at least one condition
\item \textbf{H-M4:} Test specific predictions (P1: EK-FAC favors GPT-2, P2: TracIn favors BERT, P3: TRAK invariant)
\end{itemize}
